% !TeX root = ../papers/nc-p3.tex
\section{台145の曲面の商と四次中心の束}
\label{sec:parallel-rectangle-surface}
\begin{theorem}\label{thm:parallel-rectangle-surface}
$t=\pm1,\ c=1,\ b=ay,\ A=ta,\ B=1,\ C=y,\ D=t(a+y-ay)$、
$a,y\ne0,1$ とする。これは台145の水平面と、その中の一零条件・二零条件を含む。
各符号で指定始点の同型は $(a,y)$ と $(1/a,1/y)$ の同一視で尽き、
異なる符号の同型はない。従って $s=(a+1)/(a-1),\ z=(y+1)/(y-1)$ とすると無印商は
\[
 P=s^2,\quad Q=sz,\quad R=z^2,\qquad PR=Q^2,\quad (P-1)(R-1)\ne0 .
\]
二位置移動は $P,R$ の交換である。最小周期2以下は $P=R$、すなわち $a=y$ または $ay=1$ で、
このうち $ay=a+y=1$ の一類だけが周期1である。残る全点の最小周期は4である。
\end{theorem}
\begin{proof}
二つの厚い零成分を固定した後の追加の零は $a+y-1=0$ と $a+y-ay=0$ に限られる。
二面を交換するのが同時逆数変換である。その各曲線と交点は前節のものに一致する。
四枠の復元に定理\ref{thm:third-rectangle-orbits}の部分traceと射影子を同じ順序で使う。
$\mathbb Q[a,y,1/(ay(a-1)(y-1))]$ で、位置1の文字対 $(2,3)$ の差が生成するイデアルは
$(a+1,y+1)$、他の対に用いる差は全て単位イデアルである。従って
$(-1,-1)$ の外で全四枠が復元できる。この一点は位置1で同時固有直線にならない二文字を持つので、
指定始点で他の点と同型にならない。

追加の零がないときは各始点移動で三位置分の置換を列挙し、総色から決まる残りの置換を検査する。
偶位置移動で台を保つ置換が二つ、奇位置移動では零である。各置換について薄い語は薄い語へ移り、
その対角規格化と $B=1$ は一意に厚い係数を定める。両符号で係数比の全トーラス指標を取ると、
元の始点で恒等と同時逆数変換、二移動でこれらに $a,y$ の交換を合成したものとなる。
これらの同型は台が縮む場合にも成り立つ。実際、置換・薄い語の倍率と $B$ の再規格化に使う倍率は
この二曲線上でも単元であり、一零条件上の必要性と二零条件での追加の同型は前節で求めた。
同時逆数変換は $(s,z)\mapsto(-s,-z)$ となり、商の式が従う。

固定点 $(-1,-1)$ でも二位置移動の同型は存在する。一位置移動は、
交代収縮の値が非零であるこの点で、位置0,1の隣接部分trace二乗の固有値を比較すると排除される。
正符号の多重集合は $\{-280,-280,-272,-192\}$ と $\{-456,-104,-232,-232\}$、
負符号は $\{-211,-211,-190,-160\}$ と $\{-313,-109,-175,-175\}$ である。
いずれも同じ共通の交代収縮の二乗で規格化できる。符号間の区別は次の反対階数から従う。
\end{proof}

\begin{proposition}\label{prop:parallel-surface-ranks-centers}
正符号の反対階数は $ay=1$ で8、$ay=-1$ で9、他で10である。
負符号では16から $a+y=0$ および
\[
 b^3-2b^2(a+y)+3b^2+3b-2(a+y)+1=0,\qquad b=ay,
\]
のうち成り立つ条件の数だけ引く。$ay=-1$ の全点が両符号で標準ASから外れる。
その補集合全体では五次までの経路束は標準階数の自由加群であり、正符号の四次中心の
文字別の核は階数 $(3,2,0,0)$ の自由加群となり全点で基底が特殊化する。
\end{proposition}
\begin{proof}
正符号では反対切断の色0,2,3がそれぞれ定階数2で、色1の行列式は
$(b-1)^2(b+1)/(ay(a-1)(y-1))$ である。色1の三次小行列式は共通因子 $b-1$ を除けば単位イデアルを生成し、
二次小行列式も単位イデアルを生成する。負符号では色2,3が可逆で、色0,1の三次小行列式がそれぞれ単位イデアルを生成する。
行列式の分子を取れば表示の条件を得る。正符号で第四文字のラベルを1だけ平行移動すると
奇二位置の交換対称性も全点で直接確認できる。

$y=-1/a$ とした関数体上の長さ4の次元は $(10,10,8,8)$ である。特殊化と捻りでこの障害を保つ。
逆に $ay+1$ も逆にした二変数環上で、五次までの全経路基底と四次中心の全核を求める。
各行列で、各段階の行と列の中から局所化の単元である成分だけを主元に選ぶ。
分母を払う際、行の内容を除く際、後退消去後の正規化の際にも単元以外では割らない。
経路の各総次数の次元は $(1,4,10,20,35,56)$、中心の自由列数は各文字で上記の値になる。
従って関数体の階数計算そのものがこの開集合全体の自由基底を与えている。
\path{computations/verify_parallel_rectangle_surface.py} および
\path{computations/localized_unit_echelon.py} がこれらの検算を行う。
\end{proof}
